\allowdisplaybreaks
\begin{center}
	\hrule
	\vspace{.4cm}
 \begin{tabular*}{\textwidth}{@{}l@{}|@{\extracolsep{0.6in}}r@{}}%
\parbox{4.25in}{\raggedright{\includegraphics[width=.9\linewidth]{ictp-pwf.pdf}}} &
\parbox[c][]{4in}{{\Large\textbf{Md Akiful Islam Zawad} \par}
                    { Brac University \par}
                    { KHA 224, Progati Sarani, Merul Badda, \par}
                    { Dhaka 1212, Bangladesh \par}
                    { \href{ext.akiful.islam@bracu.ac.bd}{ext.akiful.islam@bracu.ac.bd}} \par}
\end{tabular*}\vspace{.3in}
\definecolor{ceruleanblue}{rgb}{0.16, 0.32, 0.6}
	\LARGE\scshape\textbf{\textcolor{ceruleanblue}{Physics for Bangladesh QFT School}}
\end{center}
\hrule\vspace{.25in}
{\large\textbf{Problem Sheet:}\ \textsc{6} \hspace{\hfill} \large\textbf{Due Date:} \today\\
	\hrule}
 %----------------------------
\paragraph*{Problem Set 3} %\hfill \newline
\\
In this question, you are invited to explore how standard electromagnetism arises from the breaking of the electroweak force. Upon solving the previous question, you should have found that a linear combination in (\ref{eq:A-mu-def}):
\begin{align*}
    A_\mu = \cos(\theta_W) W_\mu^3 + \sin(\theta_W) B_\mu
\end{align*}
remains massless. Understand this in terms of an unbroken symmetry generated by a certain subgroup of $\text{SU(2)} \otimes \text{U(1)}$. That is, writing:
\begin{align}
    U = e^{\displaystyle i \alpha_a \frac{\sigma^a}{2} + i \beta Y},
\end{align}
(and writing $Y$ as proportional to the $2\time2$ identity matrix), identify a condition in $\alpha_a, \beta$ so that $\displaystyle\left\langle \left(i \alpha_a \frac{\sigma^a}{2} + i \beta Y\right) \varphi \right\rangle = 0$. What are the charges of the fermions under electromagnetism in terms of $g$ and $g^\prime $?
\bigskip\bigskip\hline\hline\bigskip
\subsection*{Solution}
The electroweak force is described by the gauge group $\text{SU(2)} \otimes \text{U(1)}_Y$, where $\text{SU(2)}$ is associated with the gauge bosons $W_\mu^a$, and $\text{U(1)}_Y$ is associated with the gauge boson $B_\mu$. The Higgs field $\varphi$ is a complex scalar field in the $\text{SU(2)}$ doublet representation, and it acquires a non-zero vacuum expectation value (VEV) through spontaneous symmetry breaking (SSB).

The gauge fields combine to form the physical photon $A_\mu$, the $Z$ boson $Z_\mu$, and the charged $W^\pm$ bosons. The photon $A_\mu$ remains massless, while the $W^\pm$ and $Z_\mu$ acquire masses through the Higgs mechanism.

The gauge transformation is given by:
\begin{align}
    U = e^{\displaystyle i \alpha_a \frac{\sigma^a}{2} + i \beta Y},
\end{align}
where $\sigma^a$ are the Pauli matrices corresponding to the $\text{SU(2)}$ generators, $Y$ is the hypercharge operator associated with $\text{U(1)}_Y$, and $\alpha_a(x)$ and $\beta(x)$ are local gauge parameters.

The term $\displaystyle\alpha_a \frac{\sigma^a}{2}$ represents the $\text{SU(2)}$ part of the gauge transformation. Here, $\sigma^a$ are the Pauli matrices, and they form the basis of the $\text{SU(2)}$ algebra:
\begin{align}
    \left[\sigma^a, \sigma^b\right] = 2i \epsilon^{abc} \sigma^c
\end{align}
The Pauli matrices $\sigma^a$ act on the components of the scalar field doublet $\varphi = \begin{pmatrix} \varphi_1 \\ \varphi_2 \end{pmatrix}$ or any $\text{SU(2)}$ doublet. The parameter $\alpha_a(x)$ depends on spacetime coordinates and generates a local $\text{SU(2)}$ transformation.

The hypercharge $Y$ is a generator of the $\text{U(1)}_Y$ symmetry. The parameter $\beta(x)$ is the local transformation parameter associated with $\text{U(1)}_Y$. 

For the Higgs doublet $\varphi$, the hypercharge $\displaystyle Y = \frac{1}{2}$, meaning that $\text{U(1)}_Y$ acts by multiplying the scalar field by a phase:
\begin{align}
    e^{i \beta(x) Y} = e^{\displaystyle i\frac{\beta(x)}{2}}
\end{align}
where $\beta(x)$ is a spacetime-dependent parameter.

The gauge transformation $U$ acts on the fields. For a field $\varphi$ transforming as an $\text{SU(2)}$ doublet with hypercharge $Y$, the transformation is:
\begin{align}
    \varphi(x) &\to U(x) \varphi(x) = e^{\displaystyle i \alpha_a(x) \frac{\sigma^a}{2}} e^{\displaystyle i \beta(x) Y} \varphi(x)
\end{align}
Explicitly, for the Higgs doublet $\displaystyle\varphi = \begin{pmatrix} \varphi_1 \\ \varphi_2 \end{pmatrix}$ with hypercharge $\displaystyle Y = \frac{1}{2}$, the gauge transformation acts as:
\begin{align}
    \varphi(x) &\to e^{\displaystyle i \alpha_a(x) \frac{\sigma^a}{2}} e^{\displaystyle i \beta(x) / 2} \begin{pmatrix} \varphi_1(x) \\\\ \varphi_2(x) \end{pmatrix}
\end{align}
This represents the combined action of the $\text{SU(2)}$ and $\text{U(1)}_Y$ gauge transformations on the components of the Higgs field.

The matrix exponential $U = e^{\displaystyle i \alpha_a \frac{\sigma^a}{2} + i \beta Y}$ is computed by expanding the exponential in terms of a power series:
\begin{align}
    e^{\displaystyle i \alpha_a \frac{\sigma^a}{2} + i \beta Y} = 1 + i \left( \alpha_a \frac{\sigma^a}{2} + \beta Y \right) + \frac{1}{2!} \left( i \left( \alpha_a \frac{\sigma^a}{2} + \beta Y \right) \right)^2 + \cdots
\end{align}
Because $\sigma^a$ are $2 \times 2$ matrices, this transformation is matrix-valued, acting on fields that carry $\text{SU(2)}$ charge, such as the Higgs field or left-handed fermions.

In the Higgs mechanism, the scalar field $\varphi$ has the form:
\begin{align*}
    \varphi = \frac{1}{\sqrt{2}} \begin{pmatrix} 0 \\\\ v + h(x) \end{pmatrix}
\end{align*}
where $v$ is the vacuum expectation value of the Higgs field, the condition for an unbroken symmetry is that the VEV of the Higgs field must be invariant under certain combinations of the $\text{SU(2)}$ and $\text{U(1)}_Y$ transformations.

To identify the unbroken subgroup of $\text{SU(2)} \otimes \text{U(1)}_Y$, we need to find conditions on $\alpha_a$ and $\beta$ such that the VEV of the Higgs field remains invariant:
\begin{align}
    \left\langle \left(i \alpha_a \frac{\sigma^a}{2} + i \beta Y\right) \varphi \right\rangle = 0
\end{align}
We compute the action of the gauge generator $\displaystyle i \alpha_a \frac{\sigma^a}{2} + i \beta Y$ on the Higgs VEV:
\begin{align}
    \left( \alpha_a \frac{\sigma^a}{2} + \beta Y \right) \begin{pmatrix} 0 \\ v \end{pmatrix}
    &= \begin{pmatrix} \displaystyle\frac{\alpha_1}{2} - i \displaystyle\frac{\alpha_2}{2} & \displaystyle\frac{\alpha_3}{2} \\ \\
    \displaystyle\frac{\alpha_3}{2} & -\displaystyle\frac{\alpha_1}{2} + i \displaystyle\frac{\alpha_2}{2} \end{pmatrix} \begin{pmatrix} 0 \\ \\ v \end{pmatrix} \notag \\
    &= \begin{pmatrix} v \left(\displaystyle\frac{\alpha_3}{2} + \beta Y\right) \\ 0 \end{pmatrix}
\end{align}
For the VEV to be invariant under this transformation, we must have:
\begin{align*}
    \displaystyle\frac{\alpha_3}{2} + \beta Y = 0
\end{align*}
Given that $\displaystyle Y = \frac{1}{2}$ for the Higgs field, we have:
\begin{align*}
    \alpha_3 = - \beta
\end{align*}
This implies that the combination of $\text{SU(2)}$ and $\text{U(1)}_Y$ generators that leaves the Higgs VEV invariant is:
\begin{align}
    Q = \frac{\sigma^3}{2} + Y\label{eq:Q-hypercharge-def}
\end{align}
This generator $Q$ corresponds to the electric charge operator, and the unbroken symmetry is identified as $\text{U(1)}_{\text{em}}$, which describes electromagnetism.

The photon $A_\mu$ arises as the combination of the neutral $\text{SU(2)}$ and $\text{U(1)}_Y$ gauge fields that remain massless after symmetry breaking. This field is given by (\ref{eq:A-mu-def}):
\begin{align*}
    A_\mu = \cos(\theta_W) W_\mu^3 + \sin(\theta_W) B_\mu
\end{align*}
where $\theta_W$ is the Weinberg angle, and $W_\mu^3$ and $B_\mu$ are the gauge bosons associated with the neutral components of $\text{SU(2)}$ and $\text{U(1)}_Y$.

Since the photon $A_\mu$ is associated with the unbroken $\text{U(1)}_{\text{em}}$ symmetry, it remains massless.

The fermions in the electroweak theory are assigned quantum numbers under both $\text{SU(2)}$ and $\text{U(1)}_Y$. Their electric charges under electromagnetism are determined by the generator $Q$ (\ref{eq:Q-hypercharge-def}):
\begin{align*}
    Q = \frac{\sigma^3}{2} + Y
\end{align*}
For a left-handed fermion doublet $\psi = \begin{pmatrix} \psi_1 \\ \psi_2 \end{pmatrix}$, we have:
\begin{align}
    Q \psi_1 &= \left( \frac{1}{2} + Y \right) \psi_1 \\
    Q \psi_2 &= \left( -\frac{1}{2} + Y \right) \psi_2
\end{align}

For example, in the Standard Model, the left-handed leptons form an $\text{SU(2)}$ doublet with hypercharge $\displaystyle Y = -\frac{1}{2}$. The electric charges of the left-handed electron and neutrino are:
\begin{align}
    Q_{\nu_L} &= \frac{1}{2} - \frac{1}{2} = 0 \\
    Q_{e_L} &= -\frac{1}{2} - \frac{1}{2} = -1
\end{align}
Thus, the neutrino is electrically neutral, and the electron has charge $-1$.

In the electroweak theory, the gauge group is $\text{SU(2)} \otimes \text{U(1)}_Y$, where the $\text{SU(2)}$ gauge bosons are denoted by $W_\mu^a$ with coupling constant $g$. The $\text{U(1)}_Y$ gauge boson is denoted by $B_\mu$ with coupling constant $g^\prime $.

After spontaneous symmetry breaking (SSB) by the Higgs mechanism, the neutral gauge bosons $W_\mu^3$ (the neutral component of the $\text{SU(2)}$ gauge field) and $B_\mu$ (the hypercharge field from $\text{U(1)}_Y$) mix to form the photon $A_\mu$ and the $Z$ boson. The photon field $A_\mu$ is the massless gauge boson corresponding to the unbroken $\text{U(1)}_{\text{em}}$ symmetry of electromagnetism.

The mixing of $W_\mu^3$ and $B_\mu$ can be written as the definition from (\ref{eq:Z-mu-def}) and (\ref{eq:A-mu-def}):
\begin{align*}
    Z_\mu &= -\sin(\theta_W) W_\mu^3 + \cos(\theta_W) B_\mu\\
    A_\mu &= \cos(\theta_W) W_\mu^3 + \sin(\theta_W) B_\mu
\end{align*}
where $\theta_W$ is the \textbf{Weinberg angle} (or \textit{weak mixing angle}).

To relate the photon coupling constant $e$ to the couplings $g$ and $g^\prime $, we begin by recalling that the covariant derivative for a field $\psi$ transforming under both $\text{SU(2)}$ and $\text{U(1)}_Y$ is:
\begin{align}
    D_\mu \psi &= \left( \partial_\mu - i g \frac{\sigma^a}{2} W_\mu^a - i g^\prime  \frac{Y}{2} B_\mu \right) \psi
\end{align}
After symmetry breaking, the covariant derivative becomes:
\begin{align}
    D_\mu \psi &= \left( \partial_\mu - i \frac{g}{2} \sigma^+ W_\mu^+ - i \frac{g}{2} \sigma^- W_\mu^- - i \frac{g}{2} \sigma^3 W_\mu^3 - i \frac{g^\prime }{2} Y B_\mu \right) \psi
\end{align}
The photon field $A_\mu$ couples to the electric charge $Q$, which is related to the $\text{SU(2)}$ generator $\displaystyle \frac{\sigma^3}{2}$ and the hypercharge $Y$ by the equation:
\begin{align}
    Q &= \frac{\sigma^3}{2} + Y
\end{align}
Substituting the expressions for $W_\mu^3$ and $B_\mu$ in terms of $A_\mu$ and $Z_\mu$, we focus on the term involving $A_\mu$:
\begin{align}
    D_\mu \psi &\supset \left( \partial_\mu - i \left( g \sin(\theta_W) \frac{\sigma^3}{2} + g^\prime  \cos(\theta_W) Y \right) A_\mu \right) \psi
\end{align}
Since the photon $A_\mu$ couples to the electric charge $Q$, we require that the coefficient of $A_\mu$ be proportional to $eQ$, where $e$ is the electric charge coupling constant. Therefore, the combination of $g \sin(\theta_W)$ and $g^\prime  \cos(\theta_W)$ must match the electric charge operator $Q$:
\begin{align}
    e &= g \sin(\theta_W) = g^\prime  \cos(\theta_W).\label{eq:e-def-weingberg-angle}
\end{align}
This equation relates the electric charge coupling constant $e$ to the gauge couplings $g$ (from $\text{SU(2)}$) and $g^\prime $ (from $\text{U(1)}_Y$), as well as the weak mixing angle $\theta_W$.

We can derive the explicit relationship between $g$, $g^\prime $, and $e$ by solving for $e$ in terms of $g$ and $g^\prime $. From (\ref{eq:e-def-weingberg-angle}), we find:
\begin{align}
    \tan(\theta_W) &= \frac{g^\prime }{g}
\end{align}
Using this, we can also write $e$, using (\ref{eq;cos-weinberg-def}) and (\ref{eq;sin-weinberg-def}) as:
\begin{align}
    e &= \frac{g g^\prime }{\sqrt{g^2 + g^\prime ^2}}
\end{align}
Thus, the electric charge coupling constant $e$ is related to the $\text{SU(2)}$ and $\text{U(1)}_Y$ coupling constants by the equation:
\begin{align}
    e &= \frac{g g^\prime }{\sqrt{g^2 + g^\prime ^2}}
\end{align}
\bigskip\bigskip\hline\hline\bigskip